% Level 2, Week 8 -- Testing: Prove It Works (student workbook)
\chapterweek{8}{Testing: Prove It Works}{A test is a tiny program that checks your program, so bugs get caught before players find them.}

\section{What you will learn}
\begin{itemize}[leftmargin=1.2em]
  \item How to write a \textbf{test} with \texttt{assert}.
  \item How to run your tests with \textbf{pytest} and read \texttt{PASSED} and \texttt{FAILED}.
  \item How to choose \textbf{normal} and \textbf{edge} test cases that catch sneaky bugs.
  \item How to use failing tests as a map to find and fix bugs.
\end{itemize}

\section{The big idea: let the computer check your code}
Checking a game by hand is slow: run it, play, look at the score\ldots{} again and again. A \textbf{test} is a tiny function that
checks another function for you, in less than a second. It uses the word \texttt{assert}, which means \enquote{I am sure this is true}.

\begin{center}
%% alt: How a test works. Step 1: call clamp(99, 0, 10). Step 2: it returns 10. Step 3: compare with what you expect, 10. If they are the same, the test PASSED. If they are different, it FAILED.
\begin{tikzpicture}[font=\small,>={Stealth[length=3mm]},
  box/.style={draw=black,thick,rounded corners=3pt,minimum height=1.1cm,align=center}]
  \node[box,fill=kcsky!25,minimum width=2.6cm] (in) at (0,0) {\textbf{1. Call it}\\\texttt{clamp(99, 0, 10)}};
  \node[box,fill=kcyellow!35,minimum width=2.6cm] (out) at (4.2,0) {\textbf{2. It returns}\\\texttt{10}};
  \node[box,fill=kcorange!25,minimum width=2.8cm] (exp) at (8.4,0) {\textbf{3. Compare}\\with what I expect: \texttt{10}};
  \draw[->,thick] (in) -- (out);
  \draw[->,thick] (out) -- (exp);
  \node[box,fill=kcgreen!25,minimum width=3.4cm] (pass) at (6.2,-2) {\faCheck\ \textbf{same: PASSED}};
  \node[box,fill=kcvermillion!20,minimum width=3.4cm] (fail) at (10.6,-2) {\faTimes\ \textbf{different: FAILED}};
  \draw[->,thick] (exp) -- (pass);
  \draw[->,thick,dashed] (exp) -- (fail);
\end{tikzpicture}
\end{center}

\subsection{Try assert in the Shell first}
Type these two lines in the Thonny \textbf{Shell} (the bottom window), one at a time. Guess what will happen first!
\begin{pycode}[Try it in the Shell]
assert 2 + 2 == 4
assert 2 + 2 == 5, "math is broken"
\end{pycode}
The first line is true, so \textbf{nothing happens}. That is good news. The second line is false, so Python stops and shows
\texttt{AssertionError: math is broken}. The message after the comma is the one you wrote.

\subsection{A test is a function}
\begin{pycode}[A test for clamp(), which keeps a number between two limits]
from clamp_tools import clamp

def test_too_big_comes_down_to_high():
    assert clamp(99, 0, 10) == 10, "99 is too big, so it should become 10"
\end{pycode}
A test has three parts: \textbf{call} the function with values you choose, say what you \textbf{expect} with \texttt{==}, and add a
\textbf{message} for when it goes wrong. The function's name must start with \texttt{test\_}, so that pytest can find it.

\begin{tip}[Math helper: clamp and numbers below zero]
\texttt{clamp(value, 0, 10)} keeps a number between 0 and 10, like a fence on both sides. Too big? It comes down to 10.
Too small? It goes up to 0. \texttt{-3} (\enquote{minus three}) is 3 \textbf{below} zero (see the week 1 Math helper),
so it is too small: \texttt{clamp(-3, 0, 10)} gives \texttt{0}.
\end{tip}

\subsection{Running your tests}
\begin{itemize}[leftmargin=1.2em]
  \item \textbf{Way 1, the Run button:} open your test file and press \menu{Run}. The last three lines of the file start pytest for you.
  \item \textbf{Way 2, the programmer way:} press \menu{Run} once first (so Thonny is in the right folder). Then choose
        \menu{Tools > Open system shell\ldots}, type \texttt{python -m pytest starter.py -v} and press \keys{Return}
        (if it says \enquote{command not found}, use \texttt{python3}).
  \item \textbf{See \texttt{No module named 'pytest'}?} pytest is not installed yet. Ask an adult to do
        \menu{Tools > Manage packages\ldots}, search for \texttt{pytest} and click \menu{Install} (once, no account needed).
\end{itemize}
Here is what you see when one test fails. (We made a mistake on purpose: the test expects 5.)
\begin{outputcode}
starter.py::test_middle_value_stays_the_same PASSED     [ 33%]
starter.py::test_too_small_goes_up_to_low FAILED        [ 66%]
starter.py::test_too_big_comes_down_to_high PASSED      [100%]
================ FAILURES ================
>       assert clamp(-3, 0, 10) == 5, "-3 is too small, so it should become 0"
E       AssertionError: -3 is too small, so it should become 0
E       assert 0 == 5
E        +  where 0 = clamp(-3, 0, 10)
starter.py:21: AssertionError
============ 1 failed, 2 passed in 0.03s ============
\end{outputcode}
(Your screen shows a few more lines, and the numbers at the end may be different. That is fine.)
Read it from the \textbf{bottom up}, in three steps:
\begin{enumerate}[leftmargin=1.4em,itemsep=0pt]
  \item The \textbf{last line} is the score: 1 failed, 2 passed.
  \item The line with \texttt{>} shows the test that failed.
  \item The \texttt{E} lines say what happened: \texttt{clamp} gave back \texttt{0}, but the test expected \texttt{5}.
\end{enumerate}
Then ask: \textbf{is the code wrong, or is my expected answer wrong?}
(Here the test is wrong: $-3$ should become 0, not 5.)

\subsection{Normal and edge cases}
A \textbf{normal case} is what usually happens. An \textbf{edge case} is a strange corner: that is where bugs love to hide.
\begin{center}
\renewcommand{\arraystretch}{1.3}
\begin{tabular}{@{}>{\raggedright\arraybackslash}p{4.3cm}>{\raggedright\arraybackslash}p{4.0cm}>{\raggedright\arraybackslash}p{6.2cm}@{}}
\toprule
\textbf{Function} & \textbf{Normal case} & \textbf{Edge cases} \\ \midrule
\texttt{clamp(value, 0, 10)} & \texttt{clamp(5, 0, 10)} gives 5 & too small, too big, exactly 0, exactly 10 \\
\texttt{best\_score(scores)} & best in the middle & empty list, one score, best is first, best is last \\
\texttt{encode(msg, shift)} & \texttt{"abc"} with shift 1 & letters near z, spaces and \texttt{!}, shift 0 \\
\bottomrule
\end{tabular}
\end{center}

\begin{exercise}[Test the Clamp]
Folder: \texttt{exercises/week08/ex-test-the-clamp}. \texttt{starter.py} is \textbf{your test file}. Read \texttt{clamp\_tools.py} first.
\begin{enumerate}[leftmargin=1.4em]
  \item Press \menu{Run}. How many tests say \texttt{PASSED}? \rule{1.5cm}{0.4pt} Why do the empty ones pass?
  \answerlines{1}
  \item TODO 1 and TODO 2: write an \texttt{assert} line in each empty test. TODO 3: add a test for a value exactly on a limit.
  \item Change one expected answer to something wrong. Run. Copy the last line of the output: \answerlines{1}
  \item Change it back. Then run \texttt{check\_my\_work.py}: it tries your tests on three \textbf{broken} clamps. Catch them all!
\end{enumerate}
\end{exercise}

\begin{exercise}[Test the Best Score]
Folder: \texttt{exercises/week08/ex-test-best-score}. \texttt{best\_score(scores)} finds the highest score. Plan your tests first:
\begin{center}
\renewcommand{\arraystretch}{1.9}
\begin{tabular}{@{}p{5.6cm}p{3.6cm}p{3.4cm}@{}}
\toprule
\textbf{I call\ldots} & \textbf{I expect\ldots} & \textbf{Normal or edge?} \\ \midrule
\texttt{best\_score([3, 9, 2])} & \texttt{9} & normal \\ \midrule
\texttt{best\_score([])} & & \\ \midrule
& & \\ \midrule
& & \\
\bottomrule
\end{tabular}
\end{center}
Write the tests in \texttt{starter.py}, press \menu{Run}, then check with \texttt{check\_my\_work.py}.
\end{exercise}

\begin{exercise}[Fix the Bugs the Tests Found]
Folder: \texttt{exercises/week08/ex-fix-the-bugs}. Four game helpers, one bug each. The tests are already written.
One helper finds the \textbf{average} score (see the Math helper below).
\begin{enumerate}[leftmargin=1.4em]
  \item Run \texttt{check\_my\_work.py}. Pick one \texttt{TRY AGAIN} and read its hint.
  \item Find that test in \texttt{test\_fix\_the\_bugs.py}. Read it like a sentence.
  \item Fix \textbf{one} function in \texttt{starter.py}. Never change the test! Run the checks again. Repeat.
\end{enumerate}
Fixed: \tickbox\ \texttt{lives\_after\_hit} \quad \tickbox\ \texttt{medal} \quad \tickbox\ \texttt{name\_tag} \quad \tickbox\ \texttt{average\_score}
\par When all four are fixed, \texttt{check\_my\_work.py} says \texttt{Amazing! All 5 checks passed.}
\end{exercise}

\begin{tip}[Math helper: the average]
The \textbf{average} is a fair share. Add up all the scores, then divide by how many scores there are.
Scores 10, 20 and 40: $10 + 20 + 40 = 70$, and there are 3 scores, so the average is $70 \div 3 = 23.33\ldots$
\texttt{round()} turns it into the nearest whole number: 23.
With \textbf{no} scores there is nothing to share, and dividing by 0 is impossible: Python crashes with \texttt{ZeroDivisionError}.
\end{tip}

\begin{tip}[Why testing matters]
In 1999 a spacecraft was lost on its way to Mars because two parts of its software used different units, and no test checked them together.
People can only trust apps and games that have been tested. Testing your game is a way of being fair to your players.
\end{tip}

\begin{vocab}
  \vocabword{test}{A small function that checks another function and says PASSED or FAILED.}
  \vocabword{assert}{\enquote{I am sure this is true.} If it is false, Python stops with \texttt{AssertionError}.}
  \vocabword{pytest}{A helper that finds every function called \texttt{test\_\ldots}, runs it and reports the results.}
  \vocabword{normal case}{An everyday input, like a list of three scores.}
  \vocabword{edge case}{A strange corner: empty, zero, the first, the last, exactly on a limit.}
  \vocabword{expected vs actual}{What the test says should happen, and what really happened.}
\end{vocab}

\begin{reflection}
Think of a game you like. Write one edge case its programmers should test. \answerlines{1}
When a test fails, what two things could be wrong? \answerlines{2}
Which was harder: writing tests or fixing bugs? Why? \answerlines{1}
\end{reflection}
